% !TeX program = pdflatex
% void_v1: A Differential-Attention Sub-1B Language Model Specialized for Code
% and Command-Line Reasoning
%
% Author:  Krish Kumar Sharma
% Version: v0.1 (design paper, pre-training)
% Date:    September 2026

\documentclass[11pt,a4paper]{article}

\usepackage[utf8]{inputenc}
\usepackage[T1]{fontenc}
\usepackage{lmodern}
\usepackage[margin=1in]{geometry}
\usepackage{amsmath,amssymb,amsthm}
\usepackage{mathtools}
\usepackage[round,authoryear]{natbib}
\usepackage{booktabs}
\usepackage{multirow}
\usepackage{array}
\usepackage{hyperref}
\usepackage{url}
\usepackage{xcolor}
\usepackage{listings}
\usepackage{graphicx}
\usepackage{enumitem}
\usepackage{caption}
\usepackage{titlesec}

\hypersetup{
    colorlinks=true,
    linkcolor=blue!60!black,
    citecolor=blue!60!black,
    urlcolor=blue!60!black
}

\newtheorem{theorem}{Theorem}
\newtheorem{proposition}{Proposition}
\newtheorem{lemma}{Lemma}
\newtheorem{corollary}{Corollary}

\theoremstyle{definition}
\newtheorem{definition}{Definition}
\newtheorem{assumption}{Assumption}

\theoremstyle{remark}
\newtheorem{remark}{Remark}

% Notation shortcuts
\newcommand{\R}{\mathbb{R}}
\newcommand{\E}{\mathbb{E}}
\newcommand{\softmax}{\mathrm{softmax}}
\newcommand{\rmsnorm}{\mathrm{RMSNorm}}
\newcommand{\rope}{\mathrm{RoPE}}
\newcommand{\swiglu}{\mathrm{SwiGLU}}
\newcommand{\silu}{\mathrm{SiLU}}
\newcommand{\attn}{\mathrm{Attn}}
\newcommand{\diffattn}{\mathrm{DiffAttn}}
\newcommand{\Var}{\mathrm{Var}}

\title{\textbf{void\_v1}: A Differential-Attention Sub-1B Language Model \\
       Specialized for Code and Command-Line Reasoning \\[0.5em]
       \large Design Paper (v0.1, pre-training)}
\author{Krish Kumar Sharma}
\date{September 2026}

\begin{document}
\maketitle

\begin{abstract}
We present the design of \textbf{void\_v1}, a 480M-parameter decoder-only
language model targeted at code generation and command-line reasoning, to be
trained from scratch on a curated 100B-token corpus using free-tier and
low-cost cloud compute. The architecture combines a deep-and-narrow scaling
regime with grouped-query attention, RoPE positional encoding with an extended
frequency base, a mixed sliding-window / full attention pattern, register
tokens, and Multi-Token Prediction heads for training-time auxiliary loss. Our
central architectural bet is the adoption of \emph{Differential Attention} at
the sub-1B parameter scale, a mechanism previously proposed by
\citet{ye2024differential} for reducing attention noise, but not yet published
in any released model at this size class. We formalise the noise-cancellation
argument as an inequality on expected off-diagonal attention mass, derive
closed-form parameter and FLOP budgets for the architecture, and compare
void\_v1 theoretically against five reference models in the 350M--1.3B range
(Pythia-410M, MobileLLM-350M, SmolLM2-360M, Qwen2.5-Coder-0.5B,
DeepSeek-Coder-1.3B). We further describe a three-stage curriculum with
teacher distillation from Qwen2.5-Coder-7B and an accompanying release plan
including a novel command-line reasoning benchmark, \textbf{void-cli-bench}.
This document is a design paper: no empirical results are yet available. All
predicted numbers are labelled as such and are conditioned on the training
recipe stated herein.
\end{abstract}

\tableofcontents
\vspace{1em}

%======================================================================
\section{Introduction}
%======================================================================

The sub-1B parameter regime for large language models (LLMs) has, in the
last twelve months, become surprisingly competitive.
Qwen2.5-Coder-0.5B~\citep{hui2024qwen25coder} reaches HumanEval pass@1 of
40.5, exceeding several 7B-class general-purpose models on code, while
SmolLM2-1.7B~\citep{allal2025smollm2} matches or beats larger models on
common-sense reasoning through aggressive over-training and data curation.
These results demonstrate that at small scale, \emph{training recipe} ---
not raw parameter count --- is the dominant driver of downstream quality.

Yet the architectures used by these models are strikingly similar: a
Llama-family decoder~\citep{touvron2023llama} with grouped-query attention
(GQA)~\citep{ainslie2023gqa}, rotary positional
embeddings~\citep{su2024roformer}, RMSNorm~\citep{zhang2019rmsnorm}, and a
SwiGLU feed-forward network~\citep{shazeer2020swiglu}. Recent research has
proposed a range of drop-in architectural improvements
(Differential Attention~\citep{ye2024differential},
Multi-Token Prediction~\citep{gloeckle2024mtp,deepseek2024v3},
register tokens~\citep{darcet2024registers},
QK-normalization~\citep{chameleon2024,dehghani2023vit22b}, the
Muon optimizer~\citep{jordan2024muon}), but no released model below 1B
parameters currently combines these techniques.

We propose \textbf{void\_v1} to fill that gap and, in the process, plant a
falsifiable research flag: \emph{Does Differential Attention improve code
generation quality at 480M parameters when trained on a code-specialized
corpus with distillation from a code-specialist teacher, relative to an
otherwise identical standard GQA baseline?}

\paragraph{Contributions.} This design paper contributes:
\begin{enumerate}[leftmargin=1.5em,itemsep=0.2em]
    \item A concrete 480M-parameter decoder architecture (\S\ref{sec:arch})
          combining seven proven-in-2024 techniques at a scale where the
          specific combination has not been published, with parameter and
          FLOP budgets derived from first principles (\S\ref{sec:budget}).
    \item A theoretical analysis of Differential Attention as a
          noise-cancellation operator, formalised as
          Theorem~\ref{thm:diffattn} on the expected off-diagonal attention
          mass (\S\ref{sec:theory}).
    \item A theoretical comparison with five reference models in the
          350M--1.3B range under identical assumptions (\S\ref{sec:compare}).
    \item A three-stage training curriculum with teacher distillation and
          an accompanying command-line reasoning benchmark
          (\S\ref{sec:training}, \S\ref{sec:eval}).
    \item An explicit risk register and ablation plan
          (\S\ref{sec:risks}) with graceful-degradation contracts for each
          novel component.
\end{enumerate}

\paragraph{Scope and limits.} This is a \emph{design paper}. No pretraining
run has yet been executed. Every quantitative statement about void\_v1 in
this document is either (a) a mathematical derivation from the stated
configuration, or (b) an \emph{expected value} conditioned on the training
recipe and derived by comparison with published models of similar scale.
We explicitly mark expected values throughout.

%======================================================================
\section{Related Work}
\label{sec:related}
%======================================================================

\subsection{Sub-1B language models}

The Pythia suite~\citep{biderman2023pythia} established the modern baseline
for small-model pretraining, with Pythia-410M trained on 300B tokens of the
Pile~\citep{gao2020pile}. Meta's MobileLLM~\citep{liu2024mobilellm}
demonstrated that at the sub-1B scale, \emph{depth} is more valuable than
\emph{width} per parameter: a 350M-parameter model with 30 layers
outperforms a comparably-sized wider model on reasoning benchmarks by
substantial margins. HuggingFace's SmolLM2 line~\citep{allal2025smollm2}
pushed the field forward by treating data quality (via FineWeb-Edu and
Cosmopedia synthetic textbook data) and over-training (11T tokens on the
1.7B model) as the primary quality levers.

\subsection{Code-specialized small models}

DeepSeek-Coder~\citep{guo2024deepseekcoder} introduced repository-level
sequence packing based on file-import topological sort, yielding
substantial gains on cross-file code completion. The Qwen2.5-Coder
family~\citep{hui2024qwen25coder} demonstrated that a 494M-parameter model
trained on 5.5T tokens of code-heavy data with fill-in-the-middle (FIM)
formatting~\citep{bavarian2022fim} can reach HumanEval pass@1 of 40.5,
setting the current small-model state of the art.

\subsection{Attention mechanism variants}

Grouped-Query Attention~\citep{ainslie2023gqa} reduces the KV-cache memory
footprint by sharing K and V across groups of Q heads, with negligible
quality loss up to sharing ratios of 8:1. Mistral-7B~\citep{jiang2023mistral}
introduced sliding-window attention, later combined with full attention in
Gemma-2~\citep{team2024gemma2} to yield an alternating pattern that
preserves long-range integration while reducing per-layer attention FLOPs.
QK-normalization~\citep{chameleon2024,dehghani2023vit22b} applies RMSNorm to
queries and keys prior to the dot product, preventing attention entropy
collapse in late training.

\textbf{Differential Attention}~\citep{ye2024differential}, proposed by
Microsoft Research in late 2024, computes two attention distributions with
independently-projected query and key matrices and returns their weighted
difference. The stated motivation is analogous to a differential amplifier in
analog circuits: common-mode noise present in both attention distributions
cancels, while informative signal that differs between the two branches
survives. Ablations at 3B parameters show consistent perplexity improvements
and reduced hallucination on question-answering benchmarks. To our knowledge,
no released model below 3B parameters currently uses this mechanism.

\subsection{Training-recipe innovations}

Multi-Token Prediction~\citep{gloeckle2024mtp} adds auxiliary heads that
predict tokens $t+2, t+3, \dots$ at training time, discarded at inference.
DeepSeek-V3~\citep{deepseek2024v3} demonstrates that MTP improves both
convergence speed and downstream generation quality. Register
tokens~\citep{darcet2024registers}, originally proposed for vision
transformers, provide dedicated "attention sinks" that reduce spurious
attention on semantically empty tokens. The Muon optimizer~\citep{jordan2024muon},
which applies orthogonal matrix updates via Newton-Schulz iteration, has
demonstrated 1.5--2$\times$ sample efficiency over AdamW on
small-model benchmarks in the nanoGPT speedrun community.

%======================================================================
\section{void\_v1 Architecture}
\label{sec:arch}
%======================================================================

\subsection{Design principles}

Three principles govern the design:

\begin{enumerate}[leftmargin=1.5em,itemsep=0.2em]
    \item \textbf{Depth over width.} At $\leq 1$B parameters, we allocate the
          parameter budget toward more layers rather than wider hidden
          dimensions, following the MobileLLM ablations.
    \item \textbf{One controlled novelty.} Exactly one component
          (Differential Attention) is a novel bet; all other choices are
          drawn from published models with reproducible ablations. This
          keeps the debugging surface small when interpreting training
          failures.
    \item \textbf{Graceful degradation.} Every novel or non-standard
          component has a documented rollback path
          (\S\ref{sec:risks}) that preserves the ability to ship a
          reduced-scope model if the component underperforms.
\end{enumerate}

\subsection{Overall configuration}
\label{sec:config}

Let $d$ denote the hidden dimension, $L$ the number of transformer layers,
$H$ the number of query heads, $H_{kv}$ the number of key/value heads (with
$H \bmod H_{kv} = 0$), $d_h = d / H$ the head dimension, $d_{ff}$ the
feed-forward inner dimension, $V$ the vocabulary size, and $S$ the maximum
sequence length. The full configuration is:

\begin{table}[h]
\centering
\small
\begin{tabular}{lr}
\toprule
\textbf{Hyperparameter} & \textbf{Value} \\
\midrule
Hidden dimension $d$                                 & 960 \\
Number of transformer layers $L$                     & 32 \\
Query heads $H$                                      & 15 \\
Key/Value heads $H_{kv}$ (GQA ratio 3:1)             & 5 \\
Head dimension $d_h$                                 & 64 \\
FFN inner dimension $d_{ff}$                         & 2560 \\
Vocabulary size $V$                                  & 49{,}152 \\
Max sequence length $S$                              & 2048 \\
RoPE frequency base $\theta$                         & 500{,}000 \\
Register tokens $R$                                  & 4 \\
Number of MTP heads (training only)                  & 2 (predicting $t+2, t+3$) \\
Tied input/output embeddings                         & Yes \\
Normalisation                                        & RMSNorm ($\epsilon = 10^{-5}$) \\
Activation (FFN)                                     & SwiGLU \\
Attention pattern                                    & See \S\ref{sec:attnpattern} \\
Attention head normalisation                         & QK-norm on both branches \\
Bias in linear layers                                & None \\
Weight init                                          & $\mathcal{N}(0, 0.02)$, residual projs $\times 1/\sqrt{2L}$ \\
Dropout                                              & 0 \\
\bottomrule
\end{tabular}
\caption{void\_v1 architectural configuration.}
\label{tab:config}
\end{table}

The total parameter count $N$ derives from the configuration as follows.
Let $T_\text{embed} = V \cdot d = 49{,}152 \cdot 960 = 4.72 \times 10^7$
(counted once due to embedding tying). Per transformer layer, the
attention module contributes
\begin{align}
T_\text{attn}
    &= \underbrace{2\, d \cdot H d_h}_{\text{2 Q projections}}
     + \underbrace{2\, d \cdot H_{kv} d_h}_{\text{2 K projections}}
     + \underbrace{d \cdot H_{kv} d_h}_{\text{1 V projection}}
     + \underbrace{H d_h \cdot d}_{\text{output projection}} \nonumber \\
    &\quad + \underbrace{d + 4 d_h}_{\text{RMSNorm, QK-norm, }\lambda} \nonumber \\
    &= 2 \cdot 960 \cdot 960 + 2 \cdot 960 \cdot 320 + 960 \cdot 320 + 960 \cdot 960
     + \text{(negligible)} \nonumber \\
    &\approx 3.17 \times 10^6.
\end{align}
Note the "2 Q projections" and "2 K projections" terms: Differential
Attention (\S\ref{sec:diffattn}) uses two independently-projected
Q and K matrices. V is shared across both branches. This adds
$d \cdot (H d_h + H_{kv} d_h) = 960 \cdot (960 + 320) \approx 1.23 \times 10^6$
parameters per layer compared to a standard GQA attention module, an
overhead of $\sim 63\%$ on the attention parameters and $\sim 8\%$ on the
whole model (since attention is only part of each layer).

The SwiGLU FFN contributes
\begin{equation}
T_\text{ffn} = 3 \cdot d \cdot d_{ff} = 3 \cdot 960 \cdot 2560 = 7.37 \times 10^6.
\end{equation}
Total per layer: $T_\text{attn} + T_\text{ffn} + T_\text{norms} \approx
10.55 \times 10^6$.

The final RMSNorm adds $d = 960$ parameters. The MTP heads add
$2 \cdot d \cdot V = 2 \cdot 960 \cdot 49{,}152 \approx 9.44 \times 10^7$
parameters, but these are (a) discarded at inference and (b) can be tied to
the main LM head to eliminate the cost entirely; we adopt full tying in
void\_v1.

Summing:
\begin{equation}
N \approx T_\text{embed} + L \cdot T_\text{layer} + O(d)
     = 4.72 \times 10^7 + 32 \cdot 1.055 \times 10^7
     \approx 4.85 \times 10^8.
\end{equation}
i.e.\ $N \approx 485$M parameters, consistent with the "$\sim 480$M" claim
in the abstract.

\subsection{Differential Grouped-Query Attention}
\label{sec:diffattn}

The core novel component. For an input sequence
$\mathbf{X} \in \R^{S \times d}$, standard scaled dot-product
attention with GQA computes
\begin{align}
Q &= \mathbf{X} W_Q,\quad W_Q \in \R^{d \times H d_h}, \\
K &= \mathbf{X} W_K,\quad W_K \in \R^{d \times H_{kv} d_h}, \\
V &= \mathbf{X} W_V,\quad W_V \in \R^{d \times H_{kv} d_h}, \\
\attn(Q, K, V) &= \softmax\!\left(\frac{Q K^\top}{\sqrt{d_h}} + M\right) V,
\end{align}
where $M$ is a causal mask and K, V are broadcast across the
$H / H_{kv}$ Q heads that share each KV head.

Differential GQA introduces a second Q and K projection:
\begin{align}
Q_1 &= \rmsnorm_\text{QK}(\rope(\mathbf{X} W_{Q_1})), \\
Q_2 &= \rmsnorm_\text{QK}(\rope(\mathbf{X} W_{Q_2})), \\
K_1 &= \rmsnorm_\text{QK}(\rope(\mathbf{X} W_{K_1})), \\
K_2 &= \rmsnorm_\text{QK}(\rope(\mathbf{X} W_{K_2})), \\
V   &= \mathbf{X} W_V.
\end{align}
The attention scores are then combined:
\begin{equation}
\diffattn(\mathbf{X})
   = \left[
       \softmax\!\left(\tfrac{Q_1 K_1^\top}{\sqrt{d_h}} + M\right)
     - \lambda(t) \cdot
       \softmax\!\left(\tfrac{Q_2 K_2^\top}{\sqrt{d_h}} + M\right)
     \right] V,
\label{eq:diffattn}
\end{equation}
where $\lambda(t) \in \R$ is a learnable, time-varying scalar per layer:
\begin{equation}
\lambda(t) = \exp(\lambda_{q_1} \cdot \lambda_{k_1})
           - \exp(\lambda_{q_2} \cdot \lambda_{k_2}) + \lambda_\text{init},
\end{equation}
with $\lambda_{q_i}, \lambda_{k_i} \in \R^{d_h}$ learnable and
$\lambda_\text{init} = 0.8 - 0.6 \exp(-0.3 \cdot (\ell - 1))$ for layer
index $\ell$, following the schedule of \citet{ye2024differential}.
This schedule ensures $\lambda(0) \approx 0.2$ in the deepest layer and
$\lambda(0) \approx 0.8$ in shallow layers at initialization, biasing
early layers toward stronger differential subtraction.

Cost. Compared to standard GQA, differential GQA doubles the Q and K
projection parameters (V unchanged) and roughly doubles the attention
score computation (two softmaxes). For a training-compute breakdown, see
\S\ref{sec:budget}.

\subsection{Mixed attention pattern}
\label{sec:attnpattern}

Not all layers use full attention. Following the Mistral/Gemma-2 design,
we adopt a mixed pattern:
\begin{itemize}[leftmargin=1.5em,itemsep=0.2em]
    \item \textbf{Layers 1--4:} full attention (early integration of
          broad context).
    \item \textbf{Layers 5--24:} alternating block of
          \{sliding-window (window=512), sliding-window, sliding-window, full\},
          repeated 5 times.
    \item \textbf{Layers 25--32:} full attention (final high-level
          integration).
\end{itemize}
This distributes 12 full-attention and 20 sliding-window layers, reducing
per-token attention FLOPs by approximately 40\% at sequence length 2048
while preserving the ability of the network to integrate long-range
context.

\subsection{Multi-Token Prediction (MTP) heads}
\label{sec:mtp}

At training time only, two auxiliary linear heads $\phi_2, \phi_3$ take the
final hidden state $h_t$ and predict tokens at positions $t+2$ and $t+3$
respectively. The training loss is
\begin{equation}
\mathcal{L}_\text{MTP}
   = \mathcal{L}_\text{CE}(h_t, x_{t+1})
   + \alpha_2\, \mathcal{L}_\text{CE}(\phi_2(h_t), x_{t+2})
   + \alpha_3\, \mathcal{L}_\text{CE}(\phi_3(h_t), x_{t+3}),
\end{equation}
with $\alpha_2 = \alpha_3 = 0.3$ following the DeepSeek-V3 recipe.
The MTP heads share weights with the main LM head via linear projections
of the hidden state through separate small transformer blocks
(one layer each), an approach that adds $\sim 2$M parameters per MTP head
and is discarded at inference.

\subsection{Register tokens}
\label{sec:registers}

Four learnable register tokens $\mathbf{r}_1, \dots, \mathbf{r}_4 \in \R^d$ are
prepended to every input sequence, making the effective input
$[\mathbf{r}_1, \dots, \mathbf{r}_4, x_1, \dots, x_S]$. These tokens receive
no supervision signal (their positions are masked in the loss) and serve
as dedicated attention sinks, reducing attention entropy on the
semantically meaningful tokens \citep{darcet2024registers}.

\subsection{Optimizer choice}
\label{sec:optim}

We use \textbf{Muon}~\citep{jordan2024muon} for hidden-layer weight
matrices (attention Q/K/V/O projections, FFN gate/up/down) and
\textbf{AdamW}~\citep{loshchilov2019adamw} for embedding, LM head, norms,
$\lambda$-scalars, and biases (of which there are none in linear layers
by design). Muon applies an orthogonal step:
\begin{equation}
W_{t+1} = W_t - \eta \cdot \mathrm{NewtonSchulz}_5(G_t),
\end{equation}
where $G_t$ is the gradient and $\mathrm{NewtonSchulz}_5$ approximates
the polar-decomposition orthogonal factor via five Newton-Schulz
iterations. Empirically, Muon achieves 1.5--2$\times$ sample efficiency
over AdamW on nanoGPT-scale benchmarks and is our primary lever for
training-token efficiency.

%======================================================================
\section{Theoretical Analysis}
\label{sec:theory}
%======================================================================

We now formalise three claims that motivate the void\_v1 architecture.

\subsection{Attention noise decomposition}

\begin{definition}[Attention noise]
Fix a query index $i$ and a key sequence of length $n$. Let
$A_i = \softmax(q_i K^\top / \sqrt{d_h})$ be the row-$i$ attention
distribution. Let $\mathcal{S} \subseteq \{1,\dots,n\}$ be the set of
\emph{semantically relevant} key positions for query $i$
(under some oracle relevance measure). Define the \textbf{off-diagonal
mass} as
\begin{equation}
\mathcal{N}(A_i) = \sum_{j \notin \mathcal{S}} A_{i,j}.
\end{equation}
$\mathcal{N}(A_i) \in [0, 1]$ measures the fraction of attention
"leaked" onto irrelevant tokens.
\end{definition}

\begin{assumption}[Common-mode noise]
\label{ass:commonmode}
For any query $i$, there exist \emph{signal} distributions
$A^{(s)}_1, A^{(s)}_2 \in \R^n$ (concentrated on $\mathcal{S}$) and a
\emph{shared noise} distribution $\mathcal{N}_c \in \R^n$
(concentrated on $\{1,\dots,n\} \setminus \mathcal{S}$)
such that the standard-attention outputs of the two branches decompose as
\begin{align}
A^{(1)}_i &= A^{(s)}_1 + \mathcal{N}_c + \epsilon_1, \\
A^{(2)}_i &= A^{(s)}_2 + \mathcal{N}_c + \epsilon_2,
\end{align}
where $\epsilon_1, \epsilon_2$ are independent zero-mean residuals with
$\E[\epsilon_i] = 0$ and $\Var(\epsilon_i) = \sigma^2$.
\end{assumption}

Assumption~\ref{ass:commonmode} is the key modeling step. It asserts that
two attention branches operating on the same input, but with different Q/K
projections, will produce (i) different attempts at capturing the true
signal, and (ii) similar noise patterns on the irrelevant-token subset
$\{1,\dots,n\} \setminus \mathcal{S}$, because that noise arises from
common input features (e.g.\ high-frequency positional artifacts, repeated
sub-tokens, embedding-space clustering).

\subsection{Differential attention noise reduction}

\begin{theorem}[Off-diagonal mass reduction]
\label{thm:diffattn}
Let $A^{\diffattn}_i = A^{(1)}_i - \lambda A^{(2)}_i$ denote the
(un-renormalised) differential attention row for query $i$, with
$\lambda \in [0, 1]$. Under Assumption~\ref{ass:commonmode},
\begin{equation}
\E[\mathcal{N}(A^{\diffattn}_i)]
    = (1 - \lambda) \, \E[\mathcal{N}(\mathcal{N}_c)]
    + \E[\mathcal{N}(A^{(s)}_1 - \lambda A^{(s)}_2)]
    + O(\sigma).
\end{equation}
In particular, for $\lambda = 1$, the shared-noise contribution vanishes:
\begin{equation}
\E[\mathcal{N}(A^{\diffattn}_i)]_{\lambda=1}
    = \E[\mathcal{N}(A^{(s)}_1 - A^{(s)}_2)] + O(\sigma).
\end{equation}
\end{theorem}

\begin{proof}[Proof sketch]
By linearity of expectation and Assumption~\ref{ass:commonmode},
\begin{align}
\mathcal{N}(A^{\diffattn}_i)
    &= \sum_{j \notin \mathcal{S}}
       (A^{(1)}_{i,j} - \lambda A^{(2)}_{i,j}) \\
    &= \sum_{j \notin \mathcal{S}}
       \left[
         (A^{(s)}_{1,j} - \lambda A^{(s)}_{2,j})
         + (1 - \lambda) \mathcal{N}_{c,j}
         + (\epsilon_{1,j} - \lambda \epsilon_{2,j})
       \right].
\end{align}
Since $A^{(s)}_1, A^{(s)}_2$ are concentrated on $\mathcal{S}$, the first
term contributes to $\mathcal{N}(\cdot)$ only through residual leakage,
which we absorb into $\E[\mathcal{N}(A^{(s)}_1 - \lambda A^{(s)}_2)]$.
Taking expectations, the residual term
$\E[\epsilon_{1,j} - \lambda \epsilon_{2,j}] = 0$ contributes $O(\sigma)$
(with the constant depending on $|\mathcal{S}^c|$ and $\lambda$).
The shared-noise term scales linearly with $(1 - \lambda)$, giving the
claimed decomposition.
\end{proof}

\begin{remark}
The theorem does \emph{not} claim that $A^{\diffattn}_i$ is itself a
probability distribution --- it is not, in general. Both branches
independently produce valid softmax outputs; their difference lies in
$[-1, 1]$ pointwise. In practice, the raw difference is then multiplied
by $V$ and fed through subsequent RMSNorm layers, so the value-vector
mixing recovers well-scaled activations. This is the standard treatment
in \citet{ye2024differential}. An alternative interpretation is that
$\diffattn$ produces a signed "attention field" rather than a
probability, and $\lambda$ acts as a gain control on the second branch.
\end{remark}

\begin{corollary}[Practical implication]
When Assumption~\ref{ass:commonmode} holds strongly (high common-mode
noise, low residual variance), Differential Attention with
$\lambda \to 1$ produces cleaner attention over informative tokens. When
$\lambda \to 0$, differential attention reduces to standard attention on
the first branch; the second branch is then a redundant regularizer.
The learnable schedule for $\lambda(t)$ (\S\ref{sec:diffattn}) allows
the model to interpolate between these regimes per layer and per
training step.
\end{corollary}

\subsection{Parameter efficiency of GQA}

\begin{proposition}[GQA parameter savings]
Let $\rho = H / H_{kv}$ denote the GQA sharing ratio. Compared to
multi-head attention ($\rho = 1$), grouped-query attention saves
\begin{equation}
\Delta T_\text{attn}
   = 2 \cdot d \cdot d_h \cdot H \cdot \left(1 - \frac{1}{\rho}\right)
\end{equation}
parameters per layer on the K and V projections combined. For void\_v1
($d = 960$, $d_h = 64$, $H = 15$, $\rho = 3$), this saves
$2 \cdot 960 \cdot 64 \cdot 15 \cdot (2/3) \approx 1.23 \times 10^6$
parameters per layer, or $\sim 39$M parameters across the 32-layer
model --- a $\sim 8\%$ reduction in total parameter count for
negligible quality loss per published GQA ablations \citep{ainslie2023gqa}.
Moreover, the KV cache at inference shrinks by $1 - 1/\rho = 2/3$,
reducing memory pressure during generation.
\end{proposition}

\subsection{MTP as implicit regularization}

\begin{proposition}[MTP as multi-task regularization]
The MTP training objective
\begin{equation}
\mathcal{L}_\text{MTP}
   = \mathcal{L}_\text{CE}(h_t, x_{t+1})
   + \sum_{k=2}^{K} \alpha_k \mathcal{L}_\text{CE}(\phi_k(h_t), x_{t+k})
\end{equation}
can be viewed as a multi-task learning objective in which the shared
representation $h_t$ must simultaneously support prediction of the
immediate next token and of tokens further into the future. By the
standard multi-task learning analysis \citep{caruana1997multitask}, this
biases $h_t$ toward representations that capture longer-range
predictive structure rather than only local next-token statistics,
providing an implicit regularization against overfitting to short-range
patterns.
\end{proposition}

This is a heuristic argument, not a rigorous bound. The empirical
evidence from DeepSeek-V3 (which reports both faster convergence and
better downstream evals with MTP heads) supports the interpretation.

\subsection{FLOPs and memory budget}
\label{sec:budget}

\paragraph{Forward-pass FLOPs.} For one token at sequence position $t$
with sequence length $S$:
\begin{itemize}[leftmargin=1.5em,itemsep=0.2em]
    \item Attention Q/K/V projections:
          $2 \cdot d \cdot (2 H d_h + 2 H_{kv} d_h + H_{kv} d_h) \approx
          6.34 \times 10^6$ FLOPs (dominated by diff-attn's 2$\times$ Q/K).
    \item Attention scores:
          $2 \cdot S \cdot H d_h \cdot 2 \approx 4 S \cdot d$ FLOPs
          (two softmaxes over $S$).
    \item Attention output projection: $2 \cdot d \cdot d \approx
          1.84 \times 10^6$ FLOPs.
    \item FFN: $2 \cdot 3 \cdot d \cdot d_{ff} = 1.47 \times 10^7$ FLOPs.
    \item Per-layer total: $\approx 2.3 \times 10^7$ FLOPs per token
          at $S = 2048$.
    \item 32 layers: $\approx 7.4 \times 10^8$ FLOPs per token.
    \item Embedding + final head (tied): $2 \cdot d \cdot V \approx
          9.4 \times 10^7$ FLOPs per token.
    \item MTP heads (train only): additional $\sim 3 \times 10^7$ FLOPs.
\end{itemize}
\textbf{Total forward: $\approx 8.4 \times 10^8$ FLOPs/token (train), or
$\approx 8.1 \times 10^8$ FLOPs/token (inference, no MTP)}.

\paragraph{Backward pass.} Approximately $2 \times$ the forward pass.
Total training FLOPs per token: $\approx 2.5 \times 10^9$.

\paragraph{Training-run FLOP budget.} For 100B tokens:
$100 \times 10^9 \cdot 2.5 \times 10^9 = 2.5 \times 10^{20}$ FLOPs.
On a Kaggle TPU v3-8 at $\sim 180$ TFLOPS sustained, this is
$\sim 386$ compute-hours, or approximately 4--5 months of wall-clock at
20 hours/week free allocation. On a rented H100 at $\sim 400$ TFLOPS
sustained, this is $\sim 175$ hours, or one training run for roughly
\$350--700 at spot pricing.

\paragraph{Memory (bf16 mixed precision, single GPU).}
Model parameters (bf16): $2 \cdot 485 \times 10^6 = 0.97$ GB.
Gradients (bf16): 0.97 GB.
Optimizer states (fp32 for AdamW parameters, Muon has smaller state):
$\sim 3.5$ GB.
Activations at $S = 2048$, batch $= 8$, with gradient checkpointing at
every 4 layers: $\sim 4$ GB.
KV cache during generation (GQA, batch $= 1$, full 2048 context):
$2 \cdot L \cdot H_{kv} \cdot d_h \cdot S \cdot 2 =
2 \cdot 32 \cdot 5 \cdot 64 \cdot 2048 \cdot 2 = 168$ MB.
\textbf{Total peak training memory: $\sim 10$ GB} --- fits comfortably on
a 24 GB consumer GPU or on Kaggle T4$\times$2 (with FSDP or DDP).

%======================================================================
\section{Training Methodology}
\label{sec:training}
%======================================================================

\subsection{Data pipeline}

The pretraining corpus comprises 100B tokens drawn from six sources
with the following mixture weights:

\begin{table}[h]
\centering
\small
\begin{tabular}{lrl}
\toprule
\textbf{Source} & \textbf{Share} & \textbf{Filter} \\
\midrule
The Stack v2 dedup (Python, JS, TS, Rust, Go, C, C++, Java) & 55\% & stars $\geq 5$, permissive licenses \\
FineWeb-Edu                                                  & 20\% & educational score $\geq 3$ \\
Cosmopedia v2 (synthetic textbooks)                          & 10\% & --- \\
OpenWebMath + AutoMathText                                   & 5\%  & math reasoning \\
GitHub issues + PR discussions                               & 5\%  & code-adjacent NL \\
StackExchange (code tags) + Wikipedia (en)                   & 5\%  & Q\&A, general knowledge floor \\
\bottomrule
\end{tabular}
\caption{Pretraining data mixture (100B tokens total).}
\label{tab:datamix}
\end{table}

The pipeline applies (in order): language and quality filtering,
MinHash deduplication, contamination scrubbing against every planned
evaluation benchmark (HumanEval, MBPP, MMLU, GSM8K, ARC, HellaSwag,
plus our novel void-cli-bench), FIM reformatting on 50\% of code
samples, and repository-aware sequence packing for the Stack subset.
Tokenization uses a BPE tokenizer with 49{,}152 tokens, including FIM
sentinels and ChatML special tokens for later SFT compatibility.

\subsection{Curriculum: three-stage annealing}

Following \citet{allal2025smollm2}, we adopt a three-stage curriculum:
\begin{itemize}[leftmargin=1.5em,itemsep=0.2em]
    \item \textbf{Stage 1 (85B tokens):} Broad mix per
          Table~\ref{tab:datamix}, cosine learning rate schedule with
          peak $3 \times 10^{-4}$, linear warmup over 2000 steps.
    \item \textbf{Stage 2 (10B tokens):} Anneal data mixture toward
          code and math (Stack 60\%, Cosmopedia 20\%, OpenWebMath 15\%,
          rest 5\%). Learning rate decays via cosine to $3 \times 10^{-5}$.
    \item \textbf{Stage 3 (5B tokens):} Cosmopedia and highest-quality
          code only. Add teacher distillation objective
          (\S\ref{sec:distill}). Learning rate anneals to near zero.
\end{itemize}
This three-stage schedule is empirically worth 2--5 downstream evaluation
points at negligible compute cost.

\subsection{Knowledge distillation}
\label{sec:distill}

During Stage 3, we combine the standard cross-entropy loss with a
Kullback-Leibler divergence to the output distribution of
Qwen2.5-Coder-7B~\citep{hui2024qwen25coder} (chosen as a code-specialist
teacher rather than a general-purpose teacher):
\begin{equation}
\mathcal{L}_\text{stage3}
   = (1 - \beta) \mathcal{L}_\text{CE}(p_\text{student}, y)
   + \beta \cdot T^2 \cdot \mathrm{KL}(p_\text{teacher}^{/T}
                                         \, \| \, p_\text{student}^{/T}),
\end{equation}
with temperature $T = 2$ and distillation weight $\beta = 0.5$.
Teacher logits are computed once and cached on HuggingFace Hub as
compressed tensors ($\text{top-}k=100$ per position, saving $\sim 400\times$
storage over full logits) prior to Stage 3 training.

\subsection{Post-training}

We produce two release variants:
\begin{itemize}[leftmargin=1.5em,itemsep=0.2em]
    \item \textbf{void-v1-base}: the pretrained model, model-souped from
          the last three Stage 3 checkpoints via uniform averaging.
    \item \textbf{void-v1-instruct}: base + supervised fine-tuning on
          100k instruction samples
          (self-oss-instruct-sc2-exec-filter-50k, OpenHermes-2.5 subset,
          no\_robots) + Direct Preference
          Optimization~\citep{rafailov2023dpo} on 30k pairs of
          UltraFeedback-clean.
\end{itemize}

%======================================================================
\section{Theoretical Comparison with Reference Models}
\label{sec:compare}
%======================================================================

Table~\ref{tab:compare} places void\_v1 alongside five reference
models in the 350M--1.3B parameter range. All quality figures for reference
models are published; the void\_v1 column shows the configuration
targets and \emph{expected} evaluation numbers under the training recipe
of \S\ref{sec:training}.

\begin{table}[h]
\centering
\scriptsize
\begin{tabular}{lccccccc}
\toprule
\textbf{Feature}
    & \textbf{Pythia-410M}
    & \textbf{MobileLLM-350M}
    & \textbf{SmolLM2-360M}
    & \textbf{Qwen2.5-C-0.5B}
    & \textbf{DeepSeek-C-1.3B}
    & \textbf{void\_v1} \\
\midrule
Params (M)
    & 410 & 350 & 362 & 494 & 1300 & \textbf{485} \\
Layers
    & 24 & 30 & 32 & 24 & 24 & \textbf{32} \\
Hidden
    & 1024 & 960 & 960 & 896 & 2048 & \textbf{960} \\
GQA ratio
    & 1 (MHA) & 3 & 3 & 7 & 1 (MHA) & \textbf{3} \\
RoPE $\theta$
    & 10k & 10k & 10k & 1M & 10k & \textbf{500k} \\
Attention pattern
    & full & full & full & full & full & \textbf{mixed SW/full} \\
QK-norm
    & no & no & no & no & no & \textbf{yes} \\
Differential attention
    & no & no & no & no & no & \textbf{yes} \\
MTP heads
    & no & no & no & no & no & \textbf{yes (2)} \\
Register tokens
    & no & no & no & no & no & \textbf{yes (4)} \\
FIM
    & no & no & no & yes & yes & \textbf{yes} \\
Optimizer
    & AdamW & AdamW & AdamW & AdamW & AdamW & \textbf{Muon+AdamW} \\
Training tokens
    & 300B & 1T & 4T & 5.5T & 2T & \textbf{100B target} \\
Distillation
    & no & no & no & no & no & \textbf{stage-3 from Qwen2.5-C-7B} \\
\midrule
HumanEval (pass@1)
    & 2.4 & 8.6 & 12.8 & 40.5 & 34.8 & \textbf{25--32 (expected)} \\
MBPP (pass@1)
    & $\sim 3$ & --- & 15.6 & 40.4 & 55.6 & \textbf{25--32 (expected)} \\
MMLU
    & 25.7 & 25.7 & 24.7 & 30.5 & --- & \textbf{28--32 (expected)} \\
HellaSwag
    & 40.6 & 49.6 & 54.6 & 55.5 & --- & \textbf{50--55 (expected)} \\
\bottomrule
\end{tabular}
\caption{Theoretical comparison of void\_v1 with reference sub-1.3B
         models. Reference numbers from published model cards and
         technical reports. void\_v1 numbers are \emph{expected values}
         conditioned on the training recipe of \S\ref{sec:training};
         see \S\ref{sec:expected-range} for the derivation.}
\label{tab:compare}
\end{table}

\subsection{How expected values are derived}
\label{sec:expected-range}

The expected value ranges in Table~\ref{tab:compare} are derived as follows:

\begin{enumerate}[leftmargin=1.5em,itemsep=0.2em]
    \item \textbf{Anchor point.} Qwen2.5-Coder-0.5B (494M parameters, 5.5T
          tokens) establishes the \emph{ceiling} for the size class on
          HumanEval and MBPP at $\sim 40$. This is a $55 \times$ over-training
          margin relative to Chinchilla-optimal for 494M.
    \item \textbf{Under-training penalty.} void\_v1 targets 100B tokens
          ($\sim 20 \times$ Chinchilla-optimal for 485M parameters). By the
          Chinchilla scaling laws \citep{hoffmann2022training}, halving
          the training-token budget of a small model incurs a
          $\sim 0.02$ loss increase in cross-entropy, which corresponds
          empirically to $\sim 20$--$30\%$ HumanEval degradation. Extrapolating
          from Qwen's 40.5 to our $55 \times$ smaller token budget gives a
          na\"ive floor of $\sim 20$--$25$.
    \item \textbf{Data-quality bonus.} FineWeb-Edu score-filtered data and
          Cosmopedia synthetic textbooks demonstrate $\sim 5$--$8$-point
          improvements on HellaSwag/MMLU for comparable models
          \citep{allal2025smollm2}, and analogous gains are expected for
          code metrics from our filtered Stack v2 subset.
    \item \textbf{Distillation bonus.} Stage-3 distillation from
          Qwen2.5-Coder-7B is expected to add $\sim 3$--$6$ points on
          HumanEval based on published distillation work
          \citep{team2024gemma2}.
    \item \textbf{Differential Attention bonus (speculative).} The
          published DIFF-Transformer ablations at 3B parameters show a
          $\sim 5$--$10\%$ perplexity improvement, which historically
          translates to $\sim 1$--$3$ downstream evaluation points at
          our scale. This is our research bet: the value could plausibly
          be 0 (no help at this scale) or as high as $\sim 5$ points
          (matching published ablations at larger scale).
\end{enumerate}

Summing the anchor floor with these adjustments yields the expected
HumanEval range of $\sim 25$--$32$. We stress this is a design-time
projection; the actual outcome will be determined empirically.

%======================================================================
\section{Risks and Ablation Plan}
\label{sec:risks}
%======================================================================

Every novel or non-standard component carries a documented rollback path:

\begin{table}[h]
\centering
\small
\begin{tabular}{p{4cm} p{4cm} p{6cm}}
\toprule
\textbf{Component} & \textbf{Failure mode} & \textbf{Rollback} \\
\midrule
Differential Attention & Training diverges or degrades quality after 20B tokens & Set $W_{Q_2} = W_{K_2} = 0$ at init, $\lambda = 0$; falls back to standard GQA. \\
Multi-Token Prediction & Auxiliary loss dominates or destabilises training & Set $\alpha_2 = \alpha_3 = 0$; heads are discarded and only next-token loss remains. \\
Register tokens & No measurable benefit at 20B-token milestone & Remove; retraining continues from the checkpoint with 4 fewer positions per sequence. \\
Muon optimizer & Instability or NaN loss & Swap to pure AdamW at any checkpoint; parameter update rule change only. \\
Mixed attention pattern & Long-range integration failure on evals & Convert all sliding-window layers to full attention; adds $\sim 40\%$ FLOPs per layer. \\
Teacher distillation (Stage 3) & Teacher too different from student $\to$ destabilises anneal & Drop the KL term; Stage 3 reduces to standard cross-entropy on Cosmopedia. \\
\bottomrule
\end{tabular}
\caption{Risk register with graceful-degradation rollbacks.}
\end{table}

\paragraph{Fail-fast evaluation gates.} We commit to evaluate at 20B, 50B,
85B, 95B, and 100B tokens. If HumanEval at the 20B-token milestone is
below 8 (the SmolLM2-360M level with $50 \times$ our tokens), we halt training,
diagnose, and consider rollback before spending the remaining 80\% of
compute budget.

%======================================================================
\section{Evaluation}
\label{sec:eval}
%======================================================================

\subsection{Standard benchmarks}
We report on \texttt{bigcode-evaluation-harness} (HumanEval, HumanEval-FIM,
MBPP) and \texttt{lm-eval-harness} (MMLU, HellaSwag, ARC-Easy,
ARC-Challenge, PIQA, WinoGrande, GSM8K with chain-of-thought). All
benchmarks are contamination-checked against the training corpus.

\subsection{void-cli-bench}

We introduce a novel domain-specific evaluation: a suite of $\sim 200$
command-line reasoning tasks drawn from open-source CLI codebases (fish,
zsh, ripgrep, fzf, bat, exa) covering (a) command completion given partial
input and history, (b) shell-script bug localisation, (c) man-page-based
argument recall, and (d) pipeline construction from natural-language
specifications. Contamination checks against the CLI codebase source are
enforced at data-pipeline time.

%======================================================================
\section{Conclusion}
\label{sec:conclusion}
%======================================================================

We have specified the design of void\_v1, a 485M-parameter decoder-only
language model that combines a proven-in-2024 core (deep-narrow scaling,
GQA, RoPE, RMSNorm, SwiGLU, mixed attention, MTP, register tokens,
QK-norm, Muon optimizer) with one deliberately novel bet: Differential
Attention at the sub-1B scale. We formalised the noise-cancellation
argument for Differential Attention as a theorem on expected off-diagonal
attention mass, derived the parameter and FLOP budgets from first
principles, and placed the design in a theoretical comparison with five
reference models. The design paper commits us to a specific training
recipe (100B tokens across three curriculum stages, teacher distillation
from Qwen2.5-Coder-7B, Muon+AdamW hybrid optimization) and to a fail-fast
evaluation protocol with documented rollbacks for every novel component.

The empirical validation of every claim in this document remains open.
The next milestone is a 10M-parameter debug configuration trained on
$\sim 100$M tokens, which will validate the pipeline and provide the first
empirical read on whether Differential Attention produces cleaner
attention maps at our scale before we commit to the full run.

%======================================================================
\begin{thebibliography}{99}
\small

\bibitem[Ainslie et al.(2023)]{ainslie2023gqa}
Ainslie, J., Lee-Thorp, J., de Jong, M., Zemlyanskiy, Y., Lebr\'on, F., \&
Sanghai, S. (2023). GQA: Training Generalized Multi-Query Transformer
Models from Multi-Head Checkpoints. \emph{arXiv:2305.13245}.

\bibitem[Allal et al.(2025)]{allal2025smollm2}
Allal, L. B., et al. (2025). SmolLM2: When Smol Goes Big --- Data-Centric
Training of a Small Language Model. \emph{arXiv:2502.02737}.

\bibitem[Bavarian et al.(2022)]{bavarian2022fim}
Bavarian, M., Jun, H., Tezak, N., et al. (2022). Efficient Training of
Language Models to Fill in the Middle. \emph{arXiv:2207.14255}.

\bibitem[Biderman et al.(2023)]{biderman2023pythia}
Biderman, S., et al. (2023). Pythia: A Suite for Analyzing Large Language
Models Across Training and Scaling. \emph{ICML}.

\bibitem[Caruana(1997)]{caruana1997multitask}
Caruana, R. (1997). Multitask Learning. \emph{Machine Learning}, 28, 41--75.

\bibitem[Chameleon Team(2024)]{chameleon2024}
Chameleon Team. (2024). Chameleon: Mixed-Modal Early-Fusion Foundation
Models. \emph{arXiv:2405.09818}.

\bibitem[Darcet et al.(2024)]{darcet2024registers}
Darcet, T., Oquab, M., Mairal, J., \& Bojanowski, P. (2024). Vision
Transformers Need Registers. \emph{ICLR}.

\bibitem[DeepSeek-AI(2024)]{deepseek2024v3}
DeepSeek-AI. (2024). DeepSeek-V3 Technical Report.
\emph{arXiv:2412.19437}.

\bibitem[Dehghani et al.(2023)]{dehghani2023vit22b}
Dehghani, M., et al. (2023). Scaling Vision Transformers to 22 Billion
Parameters. \emph{ICML}.

\bibitem[Gao et al.(2020)]{gao2020pile}
Gao, L., et al. (2020). The Pile: An 800GB Dataset of Diverse Text for
Language Modeling. \emph{arXiv:2101.00027}.

\bibitem[Gloeckle et al.(2024)]{gloeckle2024mtp}
Gloeckle, F., Idrissi, B. Y., Rozi\`ere, B., et al. (2024). Better \&
Faster Large Language Models via Multi-token Prediction.
\emph{arXiv:2404.19737}.

\bibitem[Guo et al.(2024)]{guo2024deepseekcoder}
Guo, D., et al. (2024). DeepSeek-Coder: When the Large Language Model
Meets Programming --- The Rise of Code Intelligence.
\emph{arXiv:2401.14196}.

\bibitem[Hoffmann et al.(2022)]{hoffmann2022training}
Hoffmann, J., et al. (2022). Training Compute-Optimal Large Language
Models. \emph{arXiv:2203.15556}.

\bibitem[Hui et al.(2024)]{hui2024qwen25coder}
Hui, B., et al. (2024). Qwen2.5-Coder Technical Report.
\emph{arXiv:2409.12186}.

\bibitem[Jiang et al.(2023)]{jiang2023mistral}
Jiang, A. Q., et al. (2023). Mistral 7B. \emph{arXiv:2310.06825}.

\bibitem[Jordan(2024)]{jordan2024muon}
Jordan, K. (2024). Muon: An Optimizer for Hidden Layers in Neural
Networks. \url{https://kellerjordan.github.io/posts/muon/}.

\bibitem[Liu et al.(2024)]{liu2024mobilellm}
Liu, Z., Zhao, C., Iandola, F., et al. (2024). MobileLLM: Optimizing
Sub-billion Parameter Language Models for On-Device Use Cases.
\emph{arXiv:2402.14905}.

\bibitem[Loshchilov \& Hutter(2019)]{loshchilov2019adamw}
Loshchilov, I., \& Hutter, F. (2019). Decoupled Weight Decay
Regularization. \emph{ICLR}.

\bibitem[Rafailov et al.(2023)]{rafailov2023dpo}
Rafailov, R., et al. (2023). Direct Preference Optimization: Your
Language Model is Secretly a Reward Model. \emph{NeurIPS}.

\bibitem[Shazeer(2020)]{shazeer2020swiglu}
Shazeer, N. (2020). GLU Variants Improve Transformer.
\emph{arXiv:2002.05202}.

\bibitem[Su et al.(2024)]{su2024roformer}
Su, J., Lu, Y., Pan, S., Murtadha, A., Wen, B., \& Liu, Y. (2024).
RoFormer: Enhanced Transformer with Rotary Position Embedding.
\emph{Neurocomputing}, 568, 127063.

\bibitem[Team(2024)]{team2024gemma2}
Gemma Team. (2024). Gemma 2: Improving Open Language Models at a
Practical Size. \emph{arXiv:2408.00118}.

\bibitem[Touvron et al.(2023)]{touvron2023llama}
Touvron, H., et al. (2023). LLaMA: Open and Efficient Foundation
Language Models. \emph{arXiv:2302.13971}.

\bibitem[Ye et al.(2024)]{ye2024differential}
Ye, T., Dong, L., Xia, Y., et al. (2024). Differential Transformer.
\emph{arXiv:2410.05258}.

\bibitem[Zhang \& Sennrich(2019)]{zhang2019rmsnorm}
Zhang, B., \& Sennrich, R. (2019). Root Mean Square Layer Normalization.
\emph{NeurIPS}.

\end{thebibliography}

\end{document}
